The broader empirical setting is one of rapid expansion and differentiation. Figure
\ref{fig:genre_proliferation} shows the growth of major metal genre families over time. The point
is not that every local scene follows the same path. It is that the paper studies an expanding and
internally differentiating creative field rather than one isolated local anomaly. That makes it
natural to ask why some places form recognizable local scenes earlier than others.

\begin{figure}[htbp]
\centering
\includegraphics[width=0.92\textwidth]{top_genre_families_over_time.png}
\caption{Global expansion and differentiation of major metal genre families.}
\label{fig:genre_proliferation}
\end{figure}

The same pattern is visible inside the genre taxonomy itself. Figure \ref{fig:subgenre_tree}
summarizes the branching structure of metal subgenres in the data. The horizontal axis records the
first year in which each subgenre appears in the Metal Archives-derived band panel, while the
vertical layout is only a visualization device. This should be read as an illustrative label tree,
not a literal genealogy. Its role in the paper is narrower: it shows that later subgenres branch
out of earlier trunks over time, so the empirical setting is one of cumulative variety creation
rather than repeated entry into a fixed category.

\begin{figure}[htbp]
\centering
\includegraphics[width=\textwidth]{subgenre_split_tree_upward.png}
\caption{Branching tree of metal subgenres. The horizontal axis records the first year each subgenre is observed in the Metal Archives-derived band data.}
\label{fig:subgenre_tree}
\end{figure}

The scene question is also spatial. Once a genre becomes legible in a few early hubs, later scenes
appear in other cities and countries. Figure \ref{fig:black_metal_diffusion} illustrates that
pattern for black metal. It does not identify a causal diffusion process and should not be read as
a single-birthplace story. Its role is simply to show that the relevant empirical object is a
sequence of local scene formations rather than one isolated national event.

\begin{figure}[htbp]
\centering
\includegraphics[width=\textwidth]{black_metal_diffusion_snapshots.png}
\caption{Black metal diffusion snapshots in 1990 and 2000. Marker area scales with the number of active local black-metal bands in each city-year, with square-root scaling for readability.}
\label{fig:black_metal_diffusion}
\end{figure}

Taken together, these figures do only two pieces of setup work. First, heavy metal generates enough
entry and differentiation to study local niche formation as an economic outcome. Second, that
process is uneven across places and unfolds through repeated local scene formation. The empirical
analysis therefore turns to the city-by-genre panel and asks which local scene characteristics
predict the transition into operational scene status. The model section takes that narrower
threshold-crossing problem, rather than the global origin of a genre, as the object to explain.
